\subsection{Marked curves and insertion dependence}

For $n\geq3$, let $C_n(\mathbb P^1)$ be the space of ordered distinct points on the complex projective line.
The diagonal action of $\operatorname{PGL}_2(\mathbb C)$ is free.
There is a unique projective transformation taking the first three points to $0,1,\infty$.
The remaining points give the explicit quotient
\[
 C_n(\mathbb P^1)/\operatorname{PGL}_2(\mathbb C)
 \cong\mathcal M_{0,n}.
\]
Thus the two spaces have complex dimensions $n$ and $n-3$.
Their compactifications retain these dimensions.
In particular, a compactification of $C_n(\mathbb P^1)$ cannot equal $\overline{\mathcal M}_{0,n+1}$, whose dimension is $n-2$.

\begin{proposition}[Insertion integrals and Euler characteristics]
\label{thm:genus0-amplitude-bar}
Put $Z=\overline{\mathcal M}_{0,n}(\mathbb C)$ with its complex orientation and real dimension $m=2n-6$.
Let $C_1,\ldots,C_n$ be complex cochain complexes.
An insertion map is a degree-zero chain map
\[
 \mu:C_1\otimes\cdots\otimes C_n\longrightarrow
 \Omega^\bullet(Z;\mathbb C),
\]
where the tensor differential has sign $(-1)^{|x_1|+\cdots+|x_{j-1}|}$ on its $j$th summand.
For closed homogeneous inputs of total degree $m$, define
\[
 \mathcal A_\mu(x_1,\ldots,x_n)=\int_Z\mu(x_1\otimes\cdots\otimes x_n).
\]
This is a multilinear function of their cohomology classes.
If, on a fixed product of input spaces, it equals an input-independent scalar for every tuple, that scalar is zero.
Therefore an identity with an Euler characteristic of a fixed complex requires separate insertion normalization, and cannot hold for arbitrary insertions when that Euler characteristic is nonzero.
\end{proposition}

\begin{proof}
The chain identity for $\mu$ sends a tensor boundary to an exact form.
Stokes' theorem on the compact oriented manifold $Z$ makes the integral of such a form zero.
Multilinearity follows from the linearity of $\mu$ and integration.
Putting one input equal to zero forces the scalar to be zero.
For an explicit instance, take $n=3$, $Z$ a point, and $C_j=\mathbb C$ in degree zero.
The multiplication map gives $\mathcal A_\mu(a,b,c)=abc$.
The Euler characteristic of the degree-zero complex $\mathbb C$ is $1$.
The tuples $(1,1,1)$ and $(0,1,1)$ give different amplitudes for the same complex.
\end{proof}

For a chiral-chain application, $\mu$ includes the actual insertion and collision maps into forms on $Z$.
Logarithmic forms along a divisor require an integration domain and boundary prescription for which the integration functional is closed.
An Euler class integral supplies a specific characteristic form, rather than an arbitrary multilinear insertion.

